\documentclass[10pt]{article}
% Math Packages
\usepackage{amsmath, mathtools}
\usepackage{amssymb}
\usepackage{amsthm}
\usepackage{amsfonts}
\usepackage{bbm}
\usepackage{breqn}
\usepackage[margin=1in]{geometry}
\usepackage{graphicx}
\usepackage{tikz}
\usetikzlibrary{arrows.meta}
\usetikzlibrary{calc}
\usepackage{forest}
\usepackage{tikz-qtree}
\graphicspath{ {./images/} }
\usepackage{hyperref}
\usepackage[capitalize]{cleveref}
\usepackage[shortlabels]{enumitem}
\usetikzlibrary{arrows,matrix,positioning}
\usepackage{multicol}

% for the pipe symbol
\usepackage[T1]{fontenc}

% Citing theorems by name. (source: https://tex.stackexchange.com/questions/109843/cleveref-and-named-theorems)
\makeatletter
\newcommand{\ncref}[1]{\cref{#1}\mynameref{#1}{\csname r@#1\endcsname}}

\def\mynameref#1#2{%
  \begingroup
  \edef\@mytxt{#2}%
  \edef\@mytst{\expandafter\@thirdoffive\@mytxt}%
  \ifx\@mytst\empty\else
  \space(\nameref{#1})\fi
  \endgroup
}
\makeatother

% Colorful Notes
\usepackage{color} \definecolor{Red}{rgb}{1,0,0} \definecolor{Blue}{rgb}{0,0,1}
\definecolor{Purple}{rgb}{.5,0,.5} \def\red{\color{Red}} \def\blue{\color{Blue}}
\def\gray{\color{gray}} \def\purple{\color{Purple}}
\newcommand{\rnote}[1]{{\red [#1]}} % \rnote{foo} gives '[foo]' in red
\newcommand{\pnote}[1]{{\purple [#1]}} % \pnote{foo} gives '[foo]' in purple
\newcommand{\bnote}[1]{{\blue #1}} % \bnote{foo} gives 'foo' in blue
\newcommand{\gnote}[1]{{\gray #1}} % \gnote{foo} gives 'foo' in gray
\newcommand{\Max}[1]{{\purple [#1]}} % \bnote{foo} then 'foo' is blue


% Claim numbering (the counter restarts after each proof environment)
\newcounter{claimcount}
\setcounter{claimcount}{0}
\newenvironment{claim}{\refstepcounter{claimcount}\par\addvspace{\medskipamount}\noindent\textbf{Claim \arabic{claimcount}:}}{}
\usepackage{etoolbox}
\AtBeginEnvironment{proof}{\setcounter{claimcount}{0}}
\newenvironment{claimproof}{\par\addvspace{\medskipamount}\noindent\textit{Proof of Claim  \arabic{claimcount}.}}{\hfill\ensuremath{\qedsymbol} \tiny{Claim}

  \medskip}
% Add claim support to cleverref
\crefname{claimcount}{Claim}{Claims}


% Math Environments
\newtheorem{theorem}{Theorem}
\newtheorem{assumption}[theorem]{Assumption}
\newtheorem{lemma}[theorem]{Lemma}
\newtheorem{proposition}[theorem]{Proposition}
\newtheorem{corollary}[theorem]{Corollary}
\newtheorem{question}[theorem]{Question}
\theoremstyle{definition}
\newtheorem{definition}[theorem]{Definition}
\newtheorem{remark}[theorem]{Remark}
\newtheorem{example}[theorem]{Example}
\newtheorem{notation}[theorem]{Notation}
\newtheorem{problem}[theorem]{Problem}

% Matrices and Column Vectors. 
\usepackage{stackengine}
\setstackgap{L}{1.0\normalbaselineskip}
\usepackage{tabstackengine}
\setstackEOL{;}% row separator
\setstackTAB{,}% column separator
\setstacktabbedgap{1ex}% inter-column gap
\setstackgap{L}{1.5\normalbaselineskip}% inter-row baselineskip
\let\nmatrix\bracketMatrixstack  %Usage: \nmatrix{1,2,3\4,5,6}
\newcommand\cv[1]{\setstackEOL{,}\bracketMatrixstack{#1}} %usage: \cv{1,2,3}

% Custom Math Coqmmands
\newcommand{\vt}{\vskip 5mm} % vertical space
\newcommand{\fl}{\noindent\textbf} % first line
\newcommand{\Fl}{\vt\noindent\textbf} % first line with space above
\newcommand{\norm}[1]{\left\lVert#1\right\rVert} % norm
\newcommand{\pnorm}[1]{\left\lVert#1\right\rVert_p} % p-norm
\newcommand{\qnorm}[1]{\left\lVert#1\right\rVert_q} % q-norm
\newcommand{\1}[1]{\textbf{1}_{\left[#1\right]}} % indicator function 
\def\limn{\lim_{n\to\infty}} % shortcut for lim as n-> infinity
\def\sumn{\sum_{n=1}^{\infty}} % shortcut for sum from n=1 to infinity
\def\sumkn{\sum_{k=1}^{n}} % shortcut for sum from k=1 to n
\def\sumin{\sum_{i=1}^{n}} % shortcut for sum from i=1 to n
\def\SAs{\sigma\text{-algebras}} % shortcut for $\sigma$-algebras
\def\SA{\sigma\text{-algebra}} % shortcut for $\sigma$-algebra
\def\Ft{\mathcal{F}_t} % time-indexed sigma-algebra (t)
\def\Fs{\mathcal{F}_s} % time-indexed sigma-algebra (s)
\def\F{\mathcal{F}} % sigma-algebra
\def\G{\mathcal{G}} % sigma-algebra
\def\R{\mathbb{R}} % Real numbers
\def\N{\mathbb{N}} % Natural numbers
\def\Z{\mathbb{Z}} % Integers
\def\E{\mathbb{E}} % Expectation
\def\P{\mathbb{P}} % Probability
\def\Q{\mathbb{Q}} % Q probability
\def\dist{\text{dist}} %Text 'dist' for things like 'dist(x,y)'
\newcommand{\indep}{\perp \!\!\! \perp}  %independence symbol
\def\Var{\mathrm{Var}} % Variance
\def\tr{\mathrm{tr}} % trace

% Brackets and Parentheses
\def\[{\left [}
    \def\]{\right ]}
% \def\({\left (}
%   \def\){\right )}



\usepackage{color}
\definecolor{Red}{rgb}{1,0,0}
\definecolor{Blue}{rgb}{0,0,1}
\definecolor{Purple}{rgb}{.5,0,.5}
\def\red{\color{Red}}
\def\blue{\color{Blue}}
\def\gray{\color{gray}}
\def\purple{\color{Purple}}
\definecolor{RoyalBlue}{cmyk}{1, 0.50, 0, 0}
\newcommand{\dempfcolor}[1]{{\color{RoyalBlue}#1}} 
\newcommand{\demph}[1]{\dempfcolor{{\sl #1}}}

% comment exactly one of the following line to show / hide the solutions
% \newcommand{\solution}[1]{{\purple #1}} % uncomment to show the solutions
\newcommand{\solution}[1]{} % uncomment to hide the solutions



\title{Lecture Notes for Math 372: \\Elementary Probability and Statistics}
\date{Last updated: \today}
% \author{mh}

\begin{document}
\begin{center}
  \section*{Math 372: Homework 07}
  \textit{Due Thursday, March 12}
\end{center}

\textit{Instructions: Do not write your solutions on this sheet---you don't have enough space to show
your work. So use separate sheets of paper. For all problems you need to show your work.}

\begin{problem}
  % Zeilberger's notes
  The continuous random variable $X$ has probability density function
  \begin{equation*}
    f(x) = \left\{ \begin{array}{l@{\quad:\quad}l}
        1/5& 0<x<1\\
        2/5& 1<x<3\\
        0& \text{elsewhere}
      \end{array}\right.
  \end{equation*}
  Let $n$ be a positive integer. Find $\E\left[X^{n} \right]$.
\end{problem}



\begin{problem}
  Let $f(x) = 3x^{2}$ for all $x$ with $0<x<1$.
  \begin{enumerate}[(a)] 
    \item Show that $f$ is a valid pdf. 
    \item Suppose $Y$ is a random variable whose pdf is $f$. Find
    $\E\left[Y \right]$.
    \item Find the variance of $Y$.
  \end{enumerate}
\end{problem}




\Fl{Note:} In the next two problems, we'll use the indicator variable notation
$\mathbf{1}_{\left[ t \geq 0\right]}$, which means the following:
\begin{equation*}
  \mathbf{1}_{\left[ t \geq 0\right]}= \left\{ \begin{array}{l@{\quad:\quad}l}
      1 & t\geq0 \\ 0 & t<0\end{array}\right.
\end{equation*}

\begin{problem}[Waiting times 1/2]
  Suppose you station yourself at a spot on an idyllic country road and watch
  as cars pass by. Let $T$ be the waiting time, in minutes, between the first
  and second cars to pass by. Assume $T$ is an exponentially distributed
  random variable with rate $\lambda>0$. In other words, the pdf of $T$ is
  \begin{equation*}
    f_{T}(t) 
    = \lambda e^{-\lambda t}\mathbf{1}_{\left[  t \geq 0\right] } 
    = \left\{ 
      \begin{array}{l@{\quad:\quad}l} \lambda e^{-\lambda t} & t \geq 0\\ 0& t<0 \end{array}\right.
  \end{equation*}

  \begin{enumerate}[(a)]
    \item Compute $\E\left[T \right]$ using the formula $\E\left[T \right] = \int_{-\infty}^{\infty} tf_{T}(t)dt$.
    \item Suppose that $\lambda=3$ cars per minute. What is the average waiting time, in seconds, for a
    car to pass by?
    \item \label{item:1} Show that
    $\P\left[ T>s+t\mid T>s\right] = \P\left[T>t \right]$ for all $t,s>0$.
    \item Interpret the equation from Part \ref{item:1} in words.
    \textit{Hint: compare with proble 2c from homework 6.} 
  \end{enumerate}
\end{problem}

\begin{problem}[Waiting times 2/2]
  Let $T$ be the time, in minutes, between successive drink orders at a bar. Suppose that a reasonable
  probability model for $T$ is given by the pdf
  \begin{equation*}
    f_{T}(t) = Ce^{-\frac{t}{4}} \mathbf{1}_{\left[ t \geq 0\right] },
  \end{equation*}
  where is $C>0$ is some number that you'll have to determine.
  \begin{enumerate}[(a)]
    \item Find a value for $C$ such that $f_{T}(t)$ is a valid pdf. (Hint: use the fact that in order to
    be a valid pdf, the function $f_{T}$ must satisfy $\int_{-\infty}^{\infty}f_{T}(t)dt=1$).
    \item What is the probability that the time between successive orders is greater than 4 minutes?
    \item On average, how long must the bartender wait between drink orders?
  \end{enumerate}
\end{problem}

\begin{problem}
  Let $X$ denote the amount of time a book on two-hour reserve is actually
  checked out, and supposed the cdf of $X$ is
  \begin{equation*}
    F(x) = \left\{ \begin{array}{l@{\quad:\quad}l} 0 & x<0\\ \frac{x^{2}}{4}& 0 \leq x<2\\ 1 &2 \leq x\end{array}\right.
  \end{equation*}
  \begin{enumerate}[(a)]
    \item Find $\P\left[X \leq 1\right]$
    \item Find $\P\left[.5 \leq X \leq 1\right]$
    \item What is the median checkout duration of $\tilde{\mu}$? [In other
    words, solve $F(\tilde{\mu})=.5$.]
    \item Obtain the density function $f(x)$
    \item Calculate $\E\left[X \right]$
    \item Calculate $\Var(X)$.
    \item If the borrower is charged an an amount $h(X)=X^{2}$ when checkout
    duration is $X$, compute the expected charge $\E\left[h(X) \right]$.
  \end{enumerate}
\end{problem}

\begin{problem}
  % Zeilberger
  Let $T$ be the waiting time until the lightbulb spontaneoulsy burns out. Assume that $T$ is
  exponentially distributed and that the expected lifetime of the bulb is 2 years.
  \begin{enumerate}[(a)]
    \item What is the rate parameter of $T$?
    \item What is the probability density function of $T$? What is the cumulative density function?
    \item Compute $\P\left[1<T<2 \right]$.
    \item Let $\sigma$ denote the standard deviation of $T$. Compute $\sigma$.
    \item Compute $\P\left[T \leq 4\sigma \right]$
    \item Compute $\P\left[T>2 \right]$.
    \item Compute $\P\left[T>4\mid T>2 \right]$.
    \item Using several complete sentences, interpret your answers to parts (d) and (e). Is there a
    contradiction?
  \end{enumerate}
\end{problem}


\begin{problem}[Exercise 4.1.2]
  Suppose the reaction temperature $X$ (in celsius) in a certain chemical
  process has uniform distribution with $A=-5$ and $B=5$.
  \begin{enumerate}[(a)]
    \item Compute $\P\left[X<0 \right]$
    \item Compute $\P\left[-2.5<X<2.5 \right]$
    \item Compute $\P\left[-2 \leq X \leq 3\right]$
    \item For $k$ satisfying $-5<k<k+4<5$, compute $\P\left[k<X<k+4 \right]$.
  \end{enumerate}
\end{problem}


\begin{problem}
  % Zeilberger
  The lifetime of a gadget costing \$300 is exponentially distributed with mean 3
  years. The manufacturer agrees to pay a full refund to a buyer if the gadget fails during the first
  year following the purchase, a two-thirds refund if it fails between the first and second year, and a
  one-third refund if it fails between the second and third year.
  If the manufacturer sells 1000 gadgets, how much should it expect to pay in refunds?
\end{problem}


\begin{problem}[Extra credit]
  An `exploding' dice roll is a dice mechanic in various boad games whereby a player gets to roll a dice,
  and then if they rolled the maximum number on the die, they get to roll it again and add the new roll
  to the old one. This can happen more than once. Compute the expectation of an exploding 6-sided dice.
\end{problem}



\newpage
% \begin{problem}
%   In this problem, we will compute $\Var(X)$, when $X\sim \text{Pois}(\lambda)$. We will utilize the
%   important series
%   \begin{equation*}
%     e^{x} = \sum_{k=0}^{\infty}\frac{x^{k}}{k!}
%   \end{equation*}
%   \begin{enumerate}[(a)]
%     \item By computing the infinite sesries, show that $\E\left[X \right] =\lambda$. 
%     % \begin{proof}
%     %   We'll show that $\E\left[X \right] = \lambda$. To see this, observe that
%     %   \begin{align*}
%     %     \E\left[X \right] 
%     %     &= \sum_{x}xp(x) \\
%     %     &= \sum_{k=0}^{\infty} k e^{-\lambda}\frac{\lambda^{k}}{k!}\\
%     %     &=e^{-\lambda} \sum_{k=0}^{\infty}\frac{k\lambda^{k}}{k!}\\ 
%     %     &= e^{-\lambda} \left( 0+\lambda+ \lambda^{2}+\frac{\lambda^{3}}{2!}+\frac{\lambda^{4}}{3!}+\frac{\lambda^{5}}{4!}+\ldots \right)\\ 
%     %     &= e^{-\lambda} \lambda\left( 1+\lambda + \frac{\lambda^{2}}{2!}+\frac{\lambda^{3}}{3!}+\ldots \right)\\ 
%     %     &= e^{-\lambda} \lambda e^{\lambda}\\
%     %     &= \lambda.
%     %   \end{align*}
%     %   The variance is a bit trickier, so we'll skip the proof of that for now.
%     % \end{proof}
%     \item Compute $\E\left[X(X-1) \right]$ using the formula $\E\left[h(X) \right] = \sum_{x}h(x)p(x)$
%     from class
%     \item Using your answer to the previous part, compute $\E\left[X^{2} \right]$. 
%     \item Using your answers to parts (a) and (c), compute $\Var(X)$.
%   \end{enumerate}
% \end{problem}

% \begin{problem}
%   Consider another approach where we show that
%   \begin{equation*}
%     \E\left[e^{tX} \right] = \exp \left( \lambda \left( e^{t}-1 \right)  \right) 
%   \end{equation*}
%   and then differentiate once or twice, and then set $t=0$. This would be a nice problem to introduce mgf.
% \end{problem}


% \begin{problem}
%   Let $X$ be a random variable, and let $\mu= \E\left[X \right]$. The variance of $X$ is defined as 
%   \begin{equation*}
%     \Var(X) = \E\left[(X-\mu)^{2} \right].
%   \end{equation*}
%   Starting from the above definition, prove that
%   \begin{equation*}
%     \Var(X) = \E\left[X^{2} \right] - \mu^{2}.
%   \end{equation*}
% \end{problem}




% \begin{problem}
%   The human genome consists of 3.2 billion nucleotides $(A,T,C,G)$. For this problem we can think of the
%   human genome as a sequence of nucleotides

%   \begin{center}
%     \texttt{CAGGAAGCATTTATATAGAGGTAATAATTAAATTTAACTTTTTATTATTCTCTATGCCCAAGAGAATGTG...}
%   \end{center}
%   and the sequence consists of 3,200,000,000 nucleotide letters.

%   The DNA mutation rate in humans is approximately $2\times 10^{-8}=0.00000002$ mutations per nucleotide
%   per generation. Assume that nucleotides mutate \textit{independently}, that is, whether one nucleotide
%   undergoes mutates is indepdenent from any other nucleotides.


%   Let $X$ be the number of mutations in
%   the genome of a randomly selected newborn baby.
%   \begin{enumerate}[(a)]
%     \item What is $\E\left[X \right]$? (Interpretation: this is the average number of mutations in a
%     human baby's genome).
%     \item Explain why $X$ is a binomial random variable. Identify the value of the two parameters $n$
%     (number of trials) and $p$ (success probability).

%     %     \item For statistical reasons that we will get to later in the class,\footnote{The Central Limit
%     %     Theorem} over 95\% of babies will have a number of mutations in the interval
%     %     \begin{equation*}
%     %       (\mu-\sigma,\mu+\sigma)
%     %     \end{equation*}
%     %     where $\mu= \E\left[X \right]$ and $\sigma = \sqrt{\Var(X)}$.
%     %     \item Compute $\Var(X)$. 
%     \item Attempt to compute some probabilities, such as $\P\left[X=1 \right]$ and $\P\left[X=60
%     \right]$, and $\P\left[X \leq 100\right]$. I do not expect you to be successful with this. What
%     makes this task difficult (even with access to a calculator or computer)?
%     %     \item Let $\lambda = \E\left[X \right]$ (which is the value you just computed). Let $Y$ be a Poisson
%     %     distribution with parameter $\lambda$. We can approximate $X$ with $Y$. Compute the following
%     %     probabilities
%     %     \begin{itemize}
%     %       \item $\P\left[Y=0 \right]$
%     %       \item $\P\left[Y=1 \right]$
%     %       \item $\P\left[Y=5 \right]$
%     %       \item $\P\left[Y \geq 3\right]$
%     %     \end{itemize}
%   \end{enumerate}
% \end{problem}



\newpage


\end{document}